% =============================================================================
% 05_report_quality.tex — Problem 5: Report Quality and Reproducibility
% =============================================================================
\section{Report Quality and Reproducibility}
\label{sec:report-quality}

\subsection{Reproduction Record}
\label{sec:reproduction}

The analysis requires R (the verified run used R 4.5.2).  Its direct
non-base dependencies, including \texttt{glmnet}, \texttt{xtable},
\texttt{jsonlite}, and \texttt{rlang}, and all transitive dependencies are
version-pinned in \texttt{renv.lock}.  The tracked \texttt{.Rprofile} points R
to the restored project library, which can be rebuilt with
\texttt{renv::restore()}.  A MiKTeX or TeX Live installation must provide
\texttt{latexmk}, XeLaTeX, and the packages listed in
\texttt{report/latex-packages.txt}; pdfLaTeX is not a supported engine because
the report contains Unicode Vietnamese names.

From the project root, the following commands restore the exact locked package
versions, run every analysis stage in dependency order, record
\texttt{sessionInfo()}, and build the report.

\begin{lstlisting}[language=bash, caption={Exact reproduction commands.}]
Rscript -e "if (!requireNamespace('renv', quietly = TRUE)) install.packages('renv', repos = 'https://cloud.r-project.org'); renv::restore(prompt = FALSE)"
Rscript R_models/00_run_all.R
cd report
latexmk -C main.tex
latexmk -xelatex -interaction=nonstopmode -halt-on-error main.tex
\end{lstlisting}

The runner executes, in order:
\texttt{01\_data\_prep\_eda.R}, \texttt{02\_ols.R},
\texttt{02\_ridge.R}, \texttt{02\_lasso.R}, \texttt{04\_enet.R},
\texttt{04\_neural.R}, \texttt{02\_comparison.R}, and
\texttt{04\_holdout.R}.  All paths are relative to the project root.

The split seed is \texttt{240201}, the shared five-fold seed is
\texttt{240301}, the fixed-feature seed is \texttt{240401}, and the bootstrap
stability seed is \texttt{240302}.  The data-preparation script creates and
stores the holdout response, but no fitting, tuning, feature-selection, or
diagnostic expression references \texttt{y\_test}.  Its first scoring use is in
\texttt{04\_holdout.R}, after the locked declaration has been written; all
candidate holdout predictions are then scored together.

The recorded software environment is included verbatim below.

\lstinputlisting[language={},basicstyle=\ttfamily\scriptsize,
  caption={R session information generated by the analysis pipeline.}]
  {../output/session_info.txt}

\subsection{Recommendation and Limitations}
\label{sec:recommendation}

\paragraph{Recommendation.}
Use the predeclared original-predictor Lasso at
$\lambda_{\min}=0.095447$ for predicting \texttt{brozek} in a comparable
adult-male population.  It had the lowest shared-fold CV RMSE (4.241) and also
the lowest one-time holdout RMSE (4.251) and MAE (3.524).  Elastic Net was a
close second on the holdout set (RMSE 4.274), so the observed advantage should
not be overstated.

\paragraph{Prediction versus interpretation.}
The recommendation prioritises prediction with a smaller deployment model:
Lasso retains nine of 14 slopes.  If coefficient stability under correlated
measurements were the primary objective, dense Ridge estimates would be less
variable in the bootstrap check, although Ridge predicts less accurately here.
Elastic Net offers a compromise.  None of these conditional coefficients or
selection events identifies causal effects or a uniquely ``true'' set of body
measurements.

\paragraph{Limitations.}
\begin{enumerate}
 \item The sample contains only 252 adult men from the source study; performance
       may not generalise to women, children, other populations, or newer
       measurement protocols.
 \item Anthropometric measurements contain measurement error, and the models
       impose predominantly linear additive structure.  The single frozen ReLU
       expansion is not an exhaustive nonlinear analysis.
 \item Tuning estimates depend on one train/test split and one shared fold
       assignment.  The 51-case holdout yields uncertain differences between
       Lasso and Elastic Net, while the random-feature result is sensitive to
       its single seed.
\end{enumerate}

\subsection{AI Verification Record}
\label{sec:ai-verification}

AI assistance is documented in Appendix~\ref{sec:ai-log}.  The group remains
responsible for explaining and independently checking every submitted item.

\begin{table}[H]
\centering
\small
\caption{Verification record for AI-assisted work.}
\label{tab:ai-verification}
\begin{tabularx}{\textwidth}{>{\raggedright\arraybackslash}p{0.20\textwidth} X X
  >{\raggedright\arraybackslash}p{0.18\textwidth}}
\toprule
\textbf{AI-Assisted Item} & \textbf{How Checked} & \textbf{Evidence} &
\textbf{Resulting Action} \\
\midrule
R pipeline and leakage audit &
Searched every script for holdout-response references and reran all eight
stages from the project root. &
Pipeline log, five unique fold IDs, generated model objects, and
\texttt{locked\_analysis.txt}. &
Kept \texttt{y\_test} scoring in \texttt{04\_holdout.R} only. \\[4pt]
Mathematical derivations &
Matched the displayed objectives to the documented \texttt{glmnet} scaling and
recomputed the eigenvalue ratio. &
$\mu_{\max}=8.942970$, $\mu_{\min}=0.002824$, and fitted condition number
15.1727. &
Corrected the Ridge $n$-scaling and normal equations. \\[4pt]
LaTeX and report audit &
Compiled with XeLaTeX, checked the log for errors/placeholders, extracted the
PDF text, and visually inspected rendered pages. &
Final PDF, build log, generated figures/tables, and Unicode author names. &
Made XeLaTeX explicit and corrected stale asset references. \\
\bottomrule
\end{tabularx}
\end{table}
